\chapter{Prescriptivism on the right and on the left}
It will be news to nobody that, in English, there has been longstanding project to remove gender-marked nouns such as \data{waitress} or \data{policeman} from the language. The project dates back to the 1960s and 1970s, when the feminist movement began to gain momentum and raise awareness about gender equality and the impact of language on societal attitudes \citep{cameron1998}. This period saw the emergence of various efforts to challenge and revise the use of gender-specific language in English \citep{lakoff1973}.

The feminist movement sought to address sexism in language by advocating for more inclusive and gender-neutral terms \citep{miller1976}. These efforts were aimed at dismantling traditional gender stereotypes and promoting a more equal representation of both men and women in various aspects of life, including the workplace, education, and government \citep{spender1980}.

Over the years, the project has made significant progress: Style guides, dictionaries, and educational institutions have increasingly adopted gender-neutral language \citep{baron1986}, and many organizations and businesses have implemented policies to promote the use of gender-inclusive terms \citep{maggio1998}. The continued efforts to make the English language more gender-neutral reflect a broader commitment to fostering a more inclusive and equitable society.

What may come as a surprise to many, as observed by Susan Neiman in her book \textit{Left is not woke}, is that in Germany, one way to avoid being sexist is precisely the opposite \citep{neiman2023}. There, at least in the context of professions and occupations, the use of gender-specific forms is considered important for promoting gender equality and inclusivity \citep{diewald2018}. The rationale behind this approach is that, by using gender-coded nouns, the German language can more accurately represent and acknowledge the presence of women in various roles, which were traditionally dominated by men. Additionally, research has shown that the use of gender-specific language in German can have a positive impact on the visibility of women and challenge traditional gender stereotypes \citep{Vervecken2012}.

The irony in this context is that, while the left is typically associated with more tolerant grammar norms and against grammatical prescriptivism, it is also much more intent on policing the use of individual words, especially slurs and words used to refer to traditionally marginalized groups. This emphasis on using inclusive language and avoiding offensive terms is seen as a way to promote equality and respect for all individuals, regardless of their background or identity. But just as there is no objectively correct way to form a subordinate interrogative in English, there is no objectively correct way to refer to individual racial groups.

As usual, that is not to say that anything goes. There are words that are clearly used with malicious intent, and such use is, rightfully, socially sanctioned. it is important to consider the complexities of language use, including in-group versus out-group dynamics and the distinction between use and mention of potentially offensive terms.

Take the example of the word \data{fucker}, which is generally considered offensive and vulgar. In some cases, friends or in-group members might use this term in a lighthearted or even affectionate manner, without intending any harm. This highlights the significance of context and intent when assessing the appropriateness of language.

Moreover, it is essential to differentiate between using a term and merely mentioning it. Discussing the word \data{fucker} in the context of this book, for instance, does not carry the same weight as directing the word at an individual in a derogatory or harmful manner.




Nunberg concludes, “racists don't use slurs because they're derogative; slurs are derogative \textit{because they're the words that racists use}.” \citep[244]{Nunberg2018}

Many disability rights theorists and organizations have rejected person-first language (e.g., \data{person with blindness}) and have advocated for identity-first language (e.g., \data{blind person}) or nominals (e.g., \data{the blind}).\footnote{This is a case of an adjective functioning as a modifier in the nominal while, at the same time functioning as its head (a Modifier–Head) in what the \textit{Cambridge grammar of the English language} called fusion of functions.} For instance the two main association working for the social and economic inequality of blind people in Canada are called The Canadian Federation of \textbf{the Blind} and the Canadian National Institute for \textbf{the Blind}. One could argue that this name is simply an anachronism to be explained away by the fact that institutions can't change their name every time linguistic preferences change, but the CFB takes the position that ``blindness is not a handicap, but a characteristic,'' and the blind and deaf communities

\section{Metphorical extensions}
A newer woke program aims to eliminate metaphorical extensions of the words when such extensions have negative connotations. ``Blindness should never be equated with darkness. And, avoid using `blind' as an adjective (e.g. \data{She was blind to the idea})'' \citep[8]{Jones2017}. The original meanings of \data{blind} in Old English include ``confused'', ``dark'', and ``obscure''. In this case, it is ``sightless'' which is the metaphorical extension, but I'll avoid the etymological fallacy. For most people today, the central meaning of \data{blind} is ``lacking sight'', and \data{being blind to something} is seen as a metaphorical use. Should we, then avoid it? One unanticipated outcome of doing so may be to deprecate use of the word in all its senses. If the only people using it are people who are bigots and the intolerant (or should that be people who are bigoted and intolerant), then the word will come to be associated only with such individuals and their attitudes. Paraphrasing Nunberg, bigots don't use these words because they're derogative; the words become derogative \data{when they're the words that bigots use}.”


In ‘Good Vibrations’ Henry Schiller talks about Beaver \& Stanley's ideas.